\begin{musica}{A Linda Rosa}{António Mafra}{Tuna de Engenharia da Universidade do Porto}
\nota{Introdução: Cantar o refrão (só vozes)}

\nota{Instrumental: G Am D7 G Am D7 G}

\ac{G}{}As rosas que te mandei
\ac{C}{}Mataram o meu desejo,
\ac{D7}{}Em troca do que te dei
\ac{G}{}Recebi de ti um beijo.
Então alegres depois
\ac{C}{}De nossos beijos trocar
\ac{D7}{}P’rá festa fomos os dois
\ac{G}{}De braço dado a cantar.

\nota{Pré-Refrão:}
\ac{G}{}Ai que bem cheira,
\ac{Am}{}Que bem cheira a linda rosa,
\ac{D7}{}Ai que bem cheira,
\ac{G}{}Que suave o seu perfume.
\ac{C}{}O – i - o - ai
\ac{G}{}Vou feliz e tu vaidosa
\ac{A7 D}{}Por uma rosa trago o coração em lume.

\nota{Refrão:}
\ac{G Am}{}Olha para a rua, o povo caminha
\ac{D G}{}Como nós vai folião.
\ac{Am}{}Minha mão na tua, tua mão na minha,
\ac{D G}{}Coração com coração.

\nota{Estrofe Instrumental}

\nota{Pré-Refrão}

\nota{Refrão}

\nota{Refrão 1 tom acima x2 – A Bm E A Bm E A}

\nota{“Ratatata” (a trêz vozes)}

\posicoes{C: x32010, D: xx0232, D7: xx0212, E: 022100, G: 320033, A: x02220,
  A7: x02020, Am: x02210, Bm: x24432}
\end{musica}
